\section{Conclusion}\label{sec:conclusion}
We built a containerised Hadoop and Spark cluster, ran it on one host with up to four
workers and on three Azure VMs, and used it to analyse 79.5~million NYC taxi trips
with three data-parallel jobs written once in Python. All engines produced identical
results. Distribution pays off only above roughly 1\,GB of input: on one month of data
the fixed start-up cost of 13--23\,s exceeds the 4.5\,s serial run. On 3.7\,GB, four
workers run 3.3--4.0$\times$ faster than one with MapReduce and 2.2--3.1$\times$ with
Spark DataFrame, and Spark DataFrame is 3.6--5.1$\times$ faster than the serial
program. The fastest engine scales worst because its fixed cost becomes the serial
fraction, so absolute times must be reported along with speedups. The combiner
reduced the shuffle by up to five orders of magnitude but saved time mainly for the
shuffle-heavy \texttt{fare} job and on the real cloud network. Two VMs scaled more
regularly than two containers on one host. A poor two-worker result on one host was
not caused by the page cache, as a controlled experiment showed, and did not reproduce
in a rerun; it shows that measurements on a shared host vary by about 20\% between
sessions.

\textbf{Future work.} (i)~Randomise and interleave the run order and repeat every
configuration in several sessions, so that host disturbances average out;
(ii)~repeat the cloud study with three or four workers in a region with free quota;
(iii)~deploy the same image on Kubernetes; (iv)~read Parquet directly in Spark.

\medskip
\noindent\textbf{Source code and video.} \TODO{Repository URL} \quad
\TODO{Demonstration video URL}
